\documentclass[11pt]{article}

\usepackage[utf8]{inputenc}
\usepackage[T1]{fontenc}
\usepackage{lmodern}
\usepackage{geometry}
\usepackage{amsmath,amssymb,amsfonts,amsthm,mathtools}
\usepackage{bm}
\usepackage{enumitem}
\usepackage{microtype}
\usepackage{hyperref}
\usepackage{titlesec}
\usepackage{booktabs}
\usepackage{array}
\usepackage{setspace}
\usepackage{mathrsfs}
\usepackage{physics}

\geometry{margin=1in}
\hypersetup{colorlinks=true, linkcolor=black, urlcolor=blue, citecolor=black}
\setlist[enumerate]{topsep=0.4em,itemsep=0.35em}
\setlist[itemize]{topsep=0.3em,itemsep=0.25em}
\titleformat{\section}{\large\bfseries}{\thesection.}{0.5em}{}
\titleformat{\subsection}{\normalsize\bfseries}{\thesubsection.}{0.5em}{}
\setstretch{1.06}

\newcommand{\Aether}{\AE{}ther}
\newcommand{\Aflow}{\AE{}ther-flow}
\newcommand{\TheoryName}{\Aflow{} Interpretation of Relativity}
\newcommand{\TheoryProgram}{\Aether{} / \Aflow{} framework}
\newcommand{\CompletedMetric}{g^{\sharp}}
\newcommand{\RootNucleus}{\mathfrak N^{\mathrm{rt}}}

\newtheorem{definition}{Definition}
\newtheorem{proposition}{Proposition}
\newtheorem{remark}{Remark}
\newtheorem{corollary}{Corollary}
\newtheorem{theorem}{Theorem}

\title{\textbf{The \TheoryName{}}\\[0.4em]
\large Primitive-Reservoir Observer-Reduction\\
\large Minimal Root Exact-Shell Transport Nucleus\\
\large Next Benchmark Criterion}
\author{}
\date{}

\begin{document}

\maketitle

\begin{abstract}
The active primitive-reservoir observer-reduction line has now moved one honest step further inside the current root layer. The previous root necessity / uniqueness theorem already showed that the benchmark-facing root core is forced once the fixed foundation package is required. The new collapse theorem then spent the remaining internal compression move available there: the larger benchmark-facing root core carries no additional benchmark-relevant freedom beyond a smaller minimal root exact-shell transport nucleus \(\RootNucleus\). After that result, another same-output root-core restatement, another root / foundation / ground relay, or a reopening of stationary-reduction, stationary-Hessian, spectral, or selector layers no longer counts as the right next benchmark-facing manuscript.

This note records the resulting criterion. The next honest benchmark-facing manuscript must therefore do one of two sharper things: either derive or justify the minimal root exact-shell transport nucleus from still more primitive substrate structure in a way that changes benchmark-facing status, or prove a genuinely smaller theorem-level collapse strictly inside \(\RootNucleus\). The benchmark boundary remains fixed throughout. The one operative metric is still exactly \(\CompletedMetric\), the ordinary matter route is unchanged, there is no second operative metric, the low-energy matter law is not enlarged, and no stronger promotion language is licensed. The positive result of the note is therefore routing discipline at the current root-nucleus frontier.
\end{abstract}

\section{Introduction}

The active derivational program of \TheoryName{} is no longer at a stage where the full benchmark-facing root core itself should be treated as the default stopping point. The previous root necessity / uniqueness theorem already removed the visible benchmark-relevant freedom at root scope. The later collapse theorem then sharpened the root situation further by proving that several shell-side constituents of that core are canonically induced rather than independently benchmark-relevant. The live frontier is therefore now the smaller minimal root exact-shell transport nucleus.

The present note records the routing consequence of that theorem. It does not reopen the root mathematics, and it does not claim a completed first-principles substrate derivation of gravity. It records only which kind of next manuscript would now count as an honest benchmark-facing gain below the current root-nucleus frontier.

\paragraph{Framework and claim boundary.}
\TheoryProgram{} denotes the broader \Aether{} / \Aflow{} framework built on the claims that the \Aether{} is the underlying four-dimensional substrate of reality, the \Aflow{} is its intrinsic ordered motion, observed three-dimensional space is the local experiential slice of that deeper substrate, S-time is the experienced order of change arising from matter, light, and the \Aflow{}, and observed expansion is the three-dimensional appearance of deeper four-dimensional ordered motion. Within that framework, \TheoryName{} denotes the exact-closure theory in which the effective gravitational dynamics are adopted to be exactly Einsteinian with universal matter coupling. In the active sequence, \emph{adoption} means use of that established relativistic dynamics without claiming substrate derivation, while \emph{derivation} is reserved for a first-principles recovery from explicit substrate variables.

The present note stays strictly inside that boundary. It records only which kind of next manuscript would now count as an honest benchmark-facing gain at the current root-nucleus frontier.

\section{Fixed Inputs}

\begin{definition}[Root-nucleus frontier fixed inputs]
The criterion recorded here uses the following fixed inputs.
\begin{enumerate}
    \item The benchmark exact-closure package, which fixes the public one-metric GR target of the repository.
    \item The post-bridge routing record and later benchmark-compression notes, which already moved the live burden upstream from same-package observer-side closure to substrate derivation of that benchmark package.
    \item The former foundation-frontier compression theorem and the root-frontier necessity / uniqueness theorem, which removed the benchmark-relevant freedom at the foundation and full-root-core stages.
    \item The root-core collapse theorem, which proves that the larger benchmark-facing root core itself collapses to the smaller minimal root exact-shell transport nucleus \(\RootNucleus\).
\end{enumerate}
\end{definition}

\begin{remark}
The present note does not replace those manuscripts. It records the routing consequence of reading them together under the active safeguard against same-output recursive depth drift.
\end{remark}

\section{Why the Root-Nucleus Collapse Theorem Is the Frontier}

\begin{definition}[Benchmark-neutral root-nucleus restatement]
A \emph{benchmark-neutral root-nucleus restatement} is a manuscript that preserves the same benchmark one-metric output, the same ordinary matter route, the same settled observer-side closure verdict, and the same minimal root exact-shell transport nucleus \(\RootNucleus\) while merely relocating the unexplained substrate layer deeper or rewriting the former larger root-core presentation without supplying a new compression, necessity, uniqueness, or comparably sharp routing gain.
\end{definition}

\begin{proposition}[The root-nucleus collapse theorem is the live frontier]
At the current root stage of the primitive-reservoir observer-reduction line, the minimal root exact-shell transport nucleus collapse theorem remains the live benchmark-facing frontier.
\end{proposition}

\begin{proof}
That theorem changes benchmark-facing status at current scope by proving a genuinely smaller theorem-level collapse strictly inside the former full root core. Once the nucleus \(\RootNucleus\) is fixed, the omitted shell-side root-core constituents are induced and carry no extra benchmark-relevant freedom. This is exactly the kind of internal compression gain allowed by the previous routing note. By contrast, a manuscript that only preserves the same benchmark package and the same root nucleus while pushing the unexplained layer deeper does not alter benchmark-facing status unless it adds a comparably sharp result. Therefore the live frontier is now the root-nucleus collapse theorem rather than the earlier larger root-core theorem or a later same-output relay.
\end{proof}

\begin{corollary}[Status of the former full-root-core theorem]
The former root necessity / uniqueness theorem remains preserved benchmark-facing main-line history, but under the current routing safeguard it is no longer by itself the default current frontier.
\end{corollary}

\begin{proof}
That theorem matters because it proved that the larger benchmark-facing root core is forced once the fixed foundation package is required. But the later collapse theorem shows that the larger core still contains redundant shell-side presentation data and therefore is not the smallest honest root-stage benchmark datum. The former theorem is preserved history and handoff, while the collapse theorem defines the current frontier.
\end{proof}

\section{The Next Benchmark Criterion}

\begin{definition}[Admissible next benchmark-facing root-nucleus manuscript]
At the current root-nucleus frontier, a manuscript counts as an admissible next benchmark-facing move only if it does at least one of the following:
\begin{enumerate}
    \item derives or justifies the minimal root exact-shell transport nucleus \(\RootNucleus\) from still more primitive substrate structure in a way that does more than merely relocate the unexplained layer deeper;
    \item proves a genuinely smaller theorem-level collapse strictly inside \(\RootNucleus\);
    \item does either of the previous two while preserving the same benchmark one-metric package, the same ordinary matter route, and the same claim boundary.
\end{enumerate}
\end{definition}

\begin{theorem}[Next honest benchmark-facing task at root-nucleus scope]
The next honest benchmark-facing manuscript at the current root-nucleus frontier must either derive or justify the minimal root exact-shell transport nucleus from still more primitive substrate structure in a benchmark-relevant way, or else prove a genuinely smaller collapse inside that nucleus. A manuscript that only repeats the former larger root core, or only repeats the same root / foundation / ground package at a deeper level without such a new gain, does not satisfy the criterion.
\end{theorem}

\begin{proof}
By Proposition 1, the live frontier already lies at the smaller root nucleus. Therefore a second manuscript below it must improve on that status rather than merely accompany it. A deeper-origin manuscript can do so only if it changes benchmark-facing status by more than depth alone. If it simply reproduces the same root nucleus through another relay, the routing rule classifies it as benchmark neutral. The remaining honest alternatives are therefore exactly those stated in the definition: either a more primitive derivation that changes benchmark-facing status, or a sharper internal collapse of the current nucleus.
\end{proof}

\section{What the Next Move Should Not Be}

\begin{proposition}[Screened next moves]
At the current root-nucleus frontier, none of the following is the default next benchmark-facing move:
\begin{enumerate}
    \item another same-output restatement of the former larger root core;
    \item another same-output root, foundation, or ground relay;
    \item a reopening of stationary-reduction, stationary-Hessian, spectral, or selector layers;
    \item a manuscript that introduces a second operative metric, enlarges the ordinary matter law, or uses stronger promotion language.
\end{enumerate}
\end{proposition}

\begin{proof}
Items (1) and (2) fail the preceding criterion because they do not directly address the current root-nucleus burden. They revisit a now-compressed larger presentation or a recorded same-output relay rather than removing the remaining nucleus-level freedom or proving a smaller internal core. Item (3) likewise does not directly address the live burden at current scope. It reopens lower or screened layers as if they were still frontier-defining rather than settling the present root-nucleus task. Item (4) violates the fixed claim boundary of the benchmark package and therefore cannot count as a disciplined next step on the active line. The benchmark package remains the exact one-metric GR package already fixed by adoption, so any admissible next manuscript must preserve that boundary.
\end{proof}

\begin{remark}
This screening statement is not a denial that those lower-layer manuscripts exist or matter historically. It is a routing statement: they are not the default way to spend the next benchmark-facing move from the current root-nucleus frontier.
\end{remark}

\section{Conclusion}

The positive result of this note is routing discipline at the current root-nucleus frontier. The collapse theorem remains the live benchmark-facing gain because it already spends the honest internal compression move still available at root scope after the larger root-core theorem: the benchmark-facing root data collapse to a smaller minimal root exact-shell transport nucleus. The former full-root-core theorem remains preserved main-line history, but under the recursive-depth safeguard it is no longer itself the default current frontier.

The scope remains narrow and honest. The benchmark package is unchanged: the one operative metric is still exactly \(\CompletedMetric\), the ordinary matter route is unchanged, there is no second operative metric, and the low-energy matter law is not enlarged. Nothing here upgrades adoption to derivation or licenses stronger promotion language.

The next open burden is therefore clear. The next benchmark-facing manuscript should derive or justify the minimal root exact-shell transport nucleus from still more primitive substrate structure in a way that changes benchmark-facing status, unless the sharper honest gain available instead is a genuinely smaller theorem-level collapse inside that nucleus. That is the clean next manuscript target at current scope.

\end{document}
